\documentclass[11pt]{article}
\usepackage[a4paper,margin=25mm]{geometry}
\usepackage{amsmath,amssymb,hyperref}
\title{Residual sampling need not improve SGD\\Conjecture 00000008776}
\author{ziangni-sys}
\date{4 October 2026}
\begin{document}
\maketitle
We disprove the explicit third-clause assertion that residual-based optimal sampling always converges faster than uniform sampling. The Chinese wording also asserts strict superiority. Identical nonzero rows are not excluded. The first clause can describe a sharp worst-case Kaczmarz bound, so it is not used here; nor are the block-size or harmonic-ratio assertions needed.

Consider the genuine consistent overdetermined system
\[
A=\begin{pmatrix}1\\1\end{pmatrix},\qquad b=\begin{pmatrix}0\\0\end{pmatrix},
\qquad x\in\mathbb R.
\]
Both rows are nonzero, and the unique solution is $x=0$. The actual individual least-squares losses and their mean are
\[
f_i(x)=\frac12(A_i x-b_i)^2=\frac{x^2}{2},\qquad
F(x)=\frac{f_1(x)+f_2(x)}2=\frac{x^2}{2}.
\]
Their gradients are all $x$, and $F$ has the unique minimizer zero. The mean normal-equation matrix is $A^TA/2=1$, so its actual diagonal inverse preconditioner is $D=1$.

At every nonzero state the two diagonally preconditioned residuals coincide. Thus sampling proportional to their magnitudes, their squares, or their positive values gives exactly probabilities $(1/2,1/2)$, the uniform law. In particular the squared-residual scores satisfy
\[
\frac{(D(A_i x-b_i))^2}{\sum_j(D(A_j x-b_j))^2}
=\frac{x^2}{2x^2}=\frac12\qquad(x\ne0).
\]
This is genuinely an optimal sampling distribution: every row gradient equals the full gradient, hence the sampling variance is zero under every row distribution and cannot be improved.

Use the safe step $\eta=1/2$ and start at $x_0=1$. For every sampled row $i$, the actual preconditioned SGD update is
\[
x-\eta D\nabla f_i(x)=x-\frac{x}{2}=\frac{x}{2}.
\]
Consequently, regardless of every random row choice,
\[
x_k=2^{-k}>0\quad\hbox{for every finite }k.
\]
The residual normalization is therefore defined at every actual finite iterate. A harmless uniform convention at zero never affects this run. Both sampling algorithms have exactly the same state distribution, namely $\delta_{2^{-k}}$, at every step. Their actual expected distance to the unique solution is exactly $2^{-k}$; their expected squared distance is $4^{-k}$; and both converge with the same factor $1/2$ in distance. Residual sampling has no strict improvement at any step or in its convergence rate.

This disproves ``always faster'' for allowed identical rows. It does not deny that a residual-based distribution can improve convergence on other systems, or that a non-strict comparison might hold under additional hypotheses.

\paragraph{Lean correspondence.}
The project defines the actual matrix and residuals, proves consistency and the genuine loss gradients/minimum, computes the mean Gram preconditioner and normalized scores, and constructs the actual uniform and residual row PMFs. SGD state laws are generated by genuine PMF bind kernels for fresh row draws, not stipulated as deterministic data. The update calculation proves that these laws are point masses for every step and every sampling kernel. Actual Bochner expectations and convergence are checked. Lean 4.19.0 and Mathlib \texttt{c44e0c8ee63ca166450922a373c7409c5d26b00b} are pinned.
\end{document}
